\documentclass[11pt]{article}
\usepackage[margin=1in]{geometry}
\usepackage[T1]{fontenc}
\usepackage{lmodern,amsmath,booktabs,graphicx}
\usepackage[hidelinks]{hyperref}
\usepackage{xurl}
\setlength{\parindent}{0pt}
\setlength{\parskip}{6pt}
\setlength{\emergencystretch}{3em}
\graphicspath{{simulation_output/estimator_visible_surface_20261003/}}
\title{Lower-Body Surface Support for Pose-Initialized LiDAR Rectangle Fitting}
\author{Pan Dynamics Collision Avoidance Research}
\date{October 3, 2026}
\begin{document}
\maketitle

\section{Implemented improvement and outcome}
The default rectangle fitter now separates edge support from containment.
Returns in the lower part of the measured body supply edge-distance and
support-midpoint residuals. Higher returns, including projected roof and trunk
surfaces, still constrain containment but are no longer forced toward footprint
edges. Position and yaw remain free variables. The NRMM prediction supplies
only their numerical initial values, with no pose penalty or locked heading.
There is one production fitting method and no new method-selection flag.

Four complete causal perception--NRMM replays were executed on two existing
300-frame recordings and their camera-blackout variants. The detector outputs
and associated native LiDAR clusters were cached results of actual YOLO26x
processing; geometry, estimator feedback and the actual MATLAB observer were
executed again. No new capture, detector inference or connection to the shared
CARLA server was made.

\begin{center}\small
\begin{tabular}{@{}llrrr@{}}\toprule
Initial gap & Metric after 2 s & Previous & Improved & Reduction \\\midrule
15 m & Raw position RMSE (m) & 0.45718 & 0.12458 & 72.75\% \\
15 m & Estimated position RMSE (m) & 0.46819 & 0.15212 & 67.51\% \\
15 m & Estimated velocity RMSE (m/s) & 0.35612 & 0.28702 & 19.40\% \\
20 m & Raw position RMSE (m) & 0.20377 & 0.04324 & 78.78\% \\
20 m & Estimated position RMSE (m) & 0.13940 & 0.05221 & 62.55\% \\
20 m & Estimated velocity RMSE (m/s) & 0.33746 & 0.11129 & 67.02\% \\
\bottomrule\end{tabular}
\end{center}
These are development-fixture results: both recordings and the same target
vehicle had already been studied. They are not an independent generalization
test. All-frame estimated position RMSE, including startup, changes from
0.40577 to 0.15444 m at 15 m and from 0.15782 to 0.10197 m at 20 m.

\section{Point selection and geometric objective}
For the current XYZ cluster $P$, define the height split
\[
 z_* = \tfrac12\{Q_5(z)+Q_{95}(z)\},\qquad
 B=\{i:z_i\leq z_*\}.
\]
If fewer than three points are selected, the threshold admits the three lowest
returns, including ties. A constant-height cluster retains every point.
This is the lower half of a robust observed \emph{height interval}, not the
lower half of the point count and not semantic surface segmentation. Its
selection uses current measurements only and is independent of the seed pose.
The percentiles and the halfway split were fixed before the first new replay;
no parameter sweep was performed.

Let $u_i=(p_{i,xy}-c)^T R(\psi)$, $h=(L/2,W/2)$ and
$e_i=|u_i|-h$. The objective is
\[
 J(c,\psi)=\frac{1}{|B|}\sum_{i\in B}
       \min\!\left(\min_j|e_{ij}|,0.5\right)^2
 +\frac{20}{|P|}\sum_{i\in P}\|\max(e_i,0)\|^2
 +\|w\odot m\|^2.
\]
The support bounds $\ell,u$ are coordinate-wise quantiles over $B$ in the
\emph{current candidate} axes, with
\[
 m=\tfrac12(\ell+u),\quad
 s=(u-\ell)\oslash(L,W),\quad
 w=\operatorname{clip}\big((s-0.8)/0.2,0,1\big).
\]
The existing continuous coverage weighting is retained. Neither the selected
support axis nor the nearest rectangle face is frozen from the prediction.
The damped finite-difference Gauss--Newton solver still optimizes all three
coordinates. Its individual step limit is a numerical safeguard, not a bound
on the final correction from the prior.

The first stage uses 2nd/98th support percentiles. The existing 0.15 m
containment-outlier rule then removes gross outliers before refinement with
minimum/maximum support. Every retained point, including upper returns,
contributes containment. Published spans and the acceptance gate now describe
the lower-body support, so roof extent does not manufacture face coverage.
The fitter records the height interval, threshold and separate boundary and
containment counts. It now requires finite XYZ input; the production pipeline
already supplies XYZ.

\begin{figure}[p]
\centering\includegraphics[width=\textwidth]{error_comparison.pdf}
\caption{Complete normal-recording error traces. The dotted line marks the
2 s reporting boundary. Startup is included in the plots and all-frame metrics.
The remaining early estimated-error peaks are distinct from the removed large
lateral geometry bias.}
\end{figure}

\section{Original maximum-error frame and remaining peaks}
At the former 15 m maximum, frame 109, time 2.18 s, simulator frame 1374:
\begin{center}\small
\begin{tabular}{@{}lrr@{}}\toprule
Quantity & Previous & Improved \\\midrule
Raw position error (m) & 0.42620 & 0.11245 \\
Raw longitudinal error (m) & 0.07213 & 0.11245 \\
Raw lateral error (m) & -0.42005 & -0.00111 \\
Fitted heading error (degrees) & -6.13715 & 1.11034 \\
Estimated position error (m) & 0.55782 & 0.19513 \\
Estimated lateral error (m) & -0.54082 & -0.00602 \\
\bottomrule\end{tabular}
\end{center}
All 999 points constrain containment; 511 also supply edge support below
$z_*=0.74543$ m. Under the new objective, the previous biased rectangle has
cost 0.02054 m$^2$, while the new fit has 0.01103 m$^2$.
The true pose with declared dimensions still has a larger cost, 0.02962 m$^2$:
the new heuristic removes the large lateral failure but does not make the
physical collision-center pose the exact minimum of this approximate model.
The remaining 0.11245 m longitudinal error is visible in the table.

At the previous raw-error maximum, frame 275, raw error falls from 0.51154
to 0.13240 m. The 20 m late estimated peak at frame 269 falls from 0.30411
to 0.04576 m. The new maxima are:
\begin{center}\small
\begin{tabular}{@{}llrrr@{}}\toprule
Gap & Window & Previous maximum (m) & New maximum (m) & New frame \\\midrule
15 m & All frames & 0.55782 & 0.33439 & 9 \\
15 m & After 2 s & 0.55782 & 0.23718 & 121 \\
20 m & All frames & 0.35801 & 0.33315 & 9 \\
20 m & After 2 s & 0.30411 & 0.09252 & 100 \\
\bottomrule\end{tabular}
\end{center}
Both all-frame maxima now occur at 0.18 s. The co-moving target-velocity
initialization remains in place, as do the original observer gains and
publication timing. At 20 m frame 9, raw error is only 0.02355 m while the
published estimate has 0.33315 m error. The preceding diagnostic investigation
isolated a 3.30156 m/s initial target-velocity mismatch in this recording;
this implementation does not correct that separate startup mechanism.

The 15 m post-2-s raw mean longitudinal error is 0.12378 m, versus lateral
mean -0.00127 m. At frame 121, raw and estimated errors are 0.12550 and
0.23718 m. Approximate vehicle dimensions and lower-body surface shape remain
plausible contributors to the longitudinal measurement bias; this experiment
does not independently identify their shares. The estimator's dynamic response
adds a transient after the raw bias changes near 1.6 s. No truth-derived
center offset or velocity has been inserted to hide either effect.

\begin{figure}[p]
\centering\includegraphics[width=\textwidth]{frame109_surface_fit.pdf}
\caption{Original worst frame: measured height separates support from
containment. Gray interior returns remain available to the containment loss.
The green reference uses the true pose and the same declared 5.0 by 1.9 m
dimensions, not a separately fitted true collision extent. The lower-body
selector greatly reduces false yaw and lateral displacement.}
\end{figure}

\section{Validation, dropout behavior and fitting cost}
All 38 Python tests pass, including six added behaviors: asymmetric dense
upper-interior returns cannot pull a fully observed footprint away from its
lower faces; upper returns still affect containment; planar rigid transforms
and height-datum shifts preserve the solution; a fully observed side can
recover pose; an unobserved face tangent remains ambiguous; invalid or absent
height is rejected. Existing tests cover free position/yaw corrections,
initial acquisition, missing observations, prefix invariance and prior timing.

The four actual MATLAB replays publish 1200 finite states. Both normal runs
have 300 valid fits. Each blackout run has 285 valid fits, with exactly frames
150--164 missing and reacquisition at 165. The normal and blackout prefixes
match exactly before interruption, ego outputs remain identical, and no
predicted pose is injected as a missing measurement. All 1170 final fit
refinements converge; 26 first-stage fits do not meet the local convergence
test before refinement. There is no claim of global optimality.

\begin{center}\small
\begin{tabular}{@{}lrrrr@{}}\toprule
Gap & \multicolumn{2}{c}{Blackout-run position RMSE after 2 s (m)}
    & \multicolumn{2}{c}{Maximum during blackout (m)} \\
 & Previous & Improved & Previous & Improved \\\midrule
15 m & 0.46580 & 0.15791 & 0.45313 & 0.22670 \\
20 m & 0.14046 & 0.05534 & 0.15531 & 0.09163 \\
\bottomrule\end{tabular}
\end{center}

An additional geometry-only comparison fits all 600 point clouds with exactly
the saved previous priors. The previous outputs reproduce within
$1.8\times10^{-15}$ m. The new fit gives post-2-s raw RMSE 0.12463 m and
0.04328 m, close to the new closed-loop values. This isolates the geometry
improvement from changes in the feedback trajectory.

That comparison alternates implementation order and times only two-stage
Python geometry. Excluding the first five calls per recording, median time
falls from 11.16 to 7.84 ms and the 95th percentile from 31.85 to 18.86 ms
over 590 timings per method. Maximum new time is 37.51 ms. Detector work,
loading and MATLAB are excluded; shared-machine timing is not an end-to-end
50 Hz guarantee. Nine selected-frame objective decompositions reproduce the
saved costs. No MATLAB unit suite was rerun because its source was unchanged;
the four actual runtime replays are the estimator integration validation.

\section{Scope, reproduction and remaining work}
The lower-body heuristic assumes observable lower surfaces and an ego
$z$ axis approximately aligned with vertical. Severe occlusion, sloped roads,
underbody contamination, different vehicle shape, density dominated by roof
returns and sparse LiDAR can violate that assumption. The synthetic rotated
side case is useful algorithm coverage, not recorded multi-view validation.
A measured boundary/shape model and causal velocity acquisition are the next
distinct improvements suggested by the remaining errors.

The native recordings retain the parked-free Town10HD\_Opt straight-follow
scenario, 143 hidden parked meshes, 0.02 s period, declared target dimensions
5.0 by 1.9 m and original LiDAR/ego-sensor seed pairs
20261002/20261003 and 20261003/20261004. Native GNSS/IMU inputs retain the
prior ideal-compass limitation. Prediction orchestration blocks opening
scoring truth; outputs are frozen before scoring. The runtime snapshot is
\path{26cd6f3998d17fd49f1cd9254245805d17f4d3ae}; the previous perception
baseline is \path{4851f4bdc3f7047cec7ad8664d3c98e8a753d0e5}.

The implementation is \path{perception/rectangle_pose_fit.py}; added tests are
\path{tests/yoloLidarSurfaceFitTest.py}. Seven compact result, audit and
provenance files accompany this report under
\path{report/ESTIMATOR_SURFACE_FIT_IMPROVEMENT_20261003/}.
Exact commands, source hashes and frozen prediction hashes are in
\path{provenance.json}. Local replay adapters, traces, test logs and figures
remain under \path{simulation_output/estimator_visible_surface_20261003/}.
They are preserved as identified external archive exports, outside the
project commit. Raw data, dependencies, generated binaries and concurrent
theory/report changes are deliberately excluded.
\end{document}
